% SOURCE_STATEMENT: "Classify $3\\times3$ Yang–Baxter operators up to Drinfeld twists, equivalence of their braid-group representations, and spectral data. What geometry do the irreducible components, dimensions, and singularities of the variety of Yang–Baxter operators have?"
% Preparation and final self-audit: CONTRIBUTING.md prompts C and D; source/audit artifacts retained in the local physics_statements_2026_10_11 batch.
% !TeX program = lualatex
\documentclass[11pt]{article}
\usepackage[margin=25mm]{geometry}
\usepackage{fontspec}
\setmainfont{Latin Modern Roman}
\usepackage{amsmath,amssymb,mathtools}
\usepackage{unicode-math}
\setmathfont{Latin Modern Math}
\usepackage{xurl,hyperref}
\hypersetup{colorlinks=true,urlcolor=blue}
\setcounter{secnumdepth}{0}
\setlength{\parindent}{0pt}
\setlength{\parskip}{5pt}
\setlength{\emergencystretch}{3em}
\newcommand{\sref}[2]{\hyperlink{source:#1:#2}{[S#2]}}
\newcommand{\cataloguescope}[1]{\paragraph{Recorded target.}#1}
\newcommand{\R}{\mathbb R}
\newcommand{\C}{\mathbb C}
\newcommand{\N}{\mathbb N}
\newcommand{\Z}{\mathbb Z}
\newcommand{\Q}{\mathbb Q}
\newcommand{\F}{\mathbb F}
\newcommand{\eps}{\varepsilon}
\begin{document}
\section{Classification and Applications of Yang–Baxter Operators}\label{statement}

\subsection{Definitions and mathematical statement}

\paragraph{Shared notation.}$\mathbb N=\{1,2,\ldots\}$; the standard number fields have their usual meanings. All additional conventions are stated locally.


A Yang--Baxter operator on $V\otimes V$ produces braid-group representations on tensor powers of $V$. Its matrix coefficients satisfy polynomial equations, defining a solution variety. Drinfeld twists provide a further equivalence relation from the source. Spectral data, irreducible components, their dimensions and singular loci are separate invariants.

\paragraph{Catalogue wording retained for traceability.}
Classify $3\times3$ Yang–Baxter operators up to Drinfeld twists, equivalence of their braid-group representations, and spectral data. What geometry do the irreducible components, dimensions, and singularities of the variety of Yang–Baxter operators have?

\paragraph{Scope and source qualification.}The source also asks whether these classifications have an impact on set-theoretic solutions; that question is retained. Its abbreviated phrase $3\times3$ operators needs a dimension convention: an operator on $V\otimes V$ for $\dim V=3$ has a $9\times9$ matrix. No silent replacement of one matrix-size convention by the other is made.
\subsection{Short English statement}
Classify the stated Yang--Baxter operators under twists and braid-representation equivalence, and determine the geometry of their solution variety.
\subsection{Sources}
\begin{itemize}
\item[\hypertarget{source:statement:1}{S1}] ULAM.AI and the UnsolvedMath Contributors, UnsolvedMath, catalogue entry. Catalogue curation under CC BY 4.0; definitions and English exposition added locally.
\url{https://huggingface.co/datasets/ulamai/UnsolvedMath}.
\item[\hypertarget{source:statement:2}{S2}] Mini-Workshop: The Yang–Baxter Equation and Representations of Braid Groups (2025), Oberwolfach Reports, DOI 10.4171/owr/2025/49, pp. 2633--2672. Original report and mathematical context. Incorporated report material remains under the source CC BY-SA 4.0 terms unless otherwise noted.
\url{https://doi.org/10.4171/owr/2025/49}.
\end{itemize}
\label{end:statement}
\end{document}
